%% ==================================================================
%%  A.R.T.E.M.I.S -- Raw Literature Review and Critical Analysis
%%
%%  This is the working (raw) version of the literature review. It is
%%  organised for internal use by the project team. A polished,
%%  paper-ready literature review will be derived from this document
%%  at a later stage.
%%
%%  Build:  pdflatex -> biber -> pdflatex -> pdflatex
%%  Backend: biblatex + biber
%% ==================================================================

\documentclass[11pt,a4paper]{report}

\usepackage[T1]{fontenc}
\usepackage[utf8]{inputenc}
\usepackage{lmodern}
\usepackage[a4paper,margin=1in]{geometry}
\usepackage{enumitem}
\usepackage{microtype}
\usepackage{titlesec}
\usepackage{csquotes}
\usepackage[hidelinks]{hyperref}

\usepackage[backend=biber,style=numeric-comp,sorting=none,giveninits=true,maxbibnames=99]{biblatex}
\addbibresource{references.bib}

%% Line spacing fixed at 1.0 (single spacing) as requested.
\linespread{1.0}

%% List formatting: tight, no dangling items, consistent labels.
\setlist[enumerate,1]{label=\arabic*., leftmargin=*, itemsep=0.3em, topsep=0.4em}
\setlist[enumerate,2]{label=\alph*., leftmargin=*, itemsep=0.2em, topsep=0.3em}
\setlist[itemize,1]{label=\textbullet, leftmargin=*, itemsep=0.3em, topsep=0.4em}
\setlist[itemize,2]{label=--, leftmargin=*, itemsep=0.2em, topsep=0.3em}

%% Numbered "part" headers within chapters, plus deep sectioning.
\setcounter{secnumdepth}{4}
\setcounter{tocdepth}{3}
\titleformat{\paragraph}[hang]{\normalfont\normalsize\itshape}{\theparagraph}{1em}{}
\titlespacing*{\paragraph}{0pt}{1.2em}{0.6em}

\newcommand{\artemis}{A.R.T.E.M.I.S}

\title{\artemis{}: Raw Literature Review and Critical Analysis\\[0.4em]
\large Working Document (Internal)}
\author{\artemis{} Project Team\\FAST-NUCES Final Year Project}
\date{September 2026}

\begin{document}

\maketitle
\tableofcontents

%% ==================================================================
\chapter{Purpose and Method of This Document}
%% ==================================================================

\section{Why This Document Exists}

This document has two jobs. The first is to record, in one place, every
paper the team read while surveying the research area around
\artemis{}, together with an accurate citation for each one. The second
is to analyse that reading critically, so that the team can state
clearly what is already known, where the open problems are, and what
each piece of work means for the design and evaluation of \artemis{}.

The document is deliberately raw. It is longer and more repetitive than
a published literature review would be, because its aim is completeness
rather than polish. A shorter, integrated literature review suitable
for inclusion in the project report will be written afterwards, using
this document as its source.

\section{How the Review Is Organised}

The reading is grouped into seven themes. Each theme corresponds to one
capability or concern of \artemis{}.

\begin{enumerate}
  \item Task interruption and resumption. This supports the
        \enquote{pick up where you left off} feature.
  \item Bounded artificial intelligence agents and trust. This supports
        the opt-in workspace model and the approval mechanism.
  \item Agentic reasoning, planning, and undo. This supports multi-step
        task execution with reversible actions.
  \item Voice interaction parity. This supports the speech input and
        output mode and its evaluation.
  \item Code assistance with human review. This supports the bounded
        code assistant feature.
  \item Hallucination in artificial intelligence summarisation. This
        concerns the note and citation synthesis feature.
  \item Retrieval-augmented generation citation accuracy. This also
        concerns the synthesis feature, specifically the correctness of
        its citations.
\end{enumerate}

\section{Structure of Each Entry}

The document has two main chapters after this one.

Chapter~\ref{ch:raw} is the raw review. For each theme it lists the
papers in turn. For every paper it records three things: what the paper
found, what research gap the paper identifies or leaves open, and what
the authors of the paper were trying to do.

Chapter~\ref{ch:analysis} is the critical analysis. For each theme it
revisits the same papers and adds a fourth element for every one: the
takeaway for \artemis{}. This states what the team should learn from
the paper, whether \artemis{} addresses the gap the paper describes,
and if so whether it does so fully or only in part.

\section{A Note on Sources}

Every source was resolved to its primary published version. Where the
initial reading recorded an archived or redirected link, the citation
here points to the peer-reviewed paper instead. Papers that existed only
as preprints, with no peer-reviewed venue, were excluded from this
review, so every entry in the reference list has a conference or journal
publication. A small number of 2026 conference and journal papers are
cited; their identifiers were checked and are genuine, not typing
errors.

%% ==================================================================
\chapter{Raw Literature Review}
\label{ch:raw}
%% ==================================================================

This chapter reports each paper on its own terms. It does not yet
connect the papers to \artemis{}; that connection is made in
Chapter~\ref{ch:analysis}.

%% ------------------------------------------------------------------
\section{Task Interruption and Resumption}
%% ------------------------------------------------------------------

Work on this problem starts from a simple observational finding: when
people are pulled away from a computer task, the presence of visible
on-screen cues, such as a window left open rather than minimised, is
what helps them pick the task back up later \autocite{iqbal2007disruption}.
This finding came out of a field study of ordinary office computer use
rather than a lab experiment, and it did not go further than
diagnosis; no system was proposed to help once leaving a window open
stops being practical, for instance after a reboot or a gap of several
days.

Later work turned this diagnosis into an actual tool. ScreenTrack
records a background time-lapse of the screen and lets a user scroll
through it visually to find and reopen something they remember
seeing, on the reasoning that people recall a sequence of visual
scenes more easily than exact file names or timestamps
\autocite{hu2020screentrack}. It performed well in a controlled lab
study for recent, visually distinct activity, but the authors were
candid about its limits: retrieval becomes slow once the history is
long, the approach depends on tasks looking visually different from
each other, so it may not suit text-heavy work, and it was never
tested in messy, long-term, multi-application use.

Taken together, these two papers describe a progression from noticing
the problem to building a partial fix that only works under fairly
narrow conditions. \artemis{}'s ``pick up where you left off'' feature
sits to one side of this line of work rather than extending it. It
does not try to imitate human visual memory or record the screen at
all; it works from the file system and from notes the user has
already written, matching on file name, recency, and explicit text
rather than on how a screen looks. Because a workspace in \artemis{}
has a bounded scope by design, the tool also avoids the ``searching a
long video'' problem that limits ScreenTrack, since there is only ever
a short list of recent, relevant items to show. This is a reasonable
substitute for the physical cue that Iqbal and Horvitz studied, but it
is worth being honest that \artemis{} has not been tested in the field
the way their study was, so the resumption benefit is currently a
design claim rather than a demonstrated one.

%% ------------------------------------------------------------------
\section{Bounded Artificial Intelligence Agents and Trust}
%% ------------------------------------------------------------------

A separate line of work asks how an AI agent's autonomy can be kept
within limits a person can actually trust. One direction is
developer-facing: a trace-based assurance framework logs everything an
agent does as a structured trace with machine-checkable contracts
attached to each step, so that a violation can be traced to an exact
point and replayed deterministically, and every action is mediated as
allow, rewrite, or block according to per-agent capability limits
\autocite{paduraru2026tracebased}. This is built to catch problems
before deployment rather than to explain a pending action to an end
user, and notably it has no way to undo an action once it has already
run.

A recent security survey looks at the same broad space from a
different angle and is more directly relevant to systems that rely on
human sign-off. It examines privilege separation as an architectural
defence, drawing a comparison to how an operating system separates a
trusted kernel from untrusted programs, and applies this to isolating
a planning component from untrusted external data
\autocite{kim2026attackdefense}. More importantly, it looks squarely
at the limits of human-in-the-loop approval itself, and its conclusion
is not encouraging: frequent approval prompts cause decision fatigue
that undercuts the safety benefit they are meant to provide, and most
people lack the background to judge what they are being asked to
approve, so they end up either approving risky actions without real
scrutiny or rejecting harmless ones. The survey names this as an open
problem for the field and does not offer a solution.

These two pieces of work show that the technical machinery for
bounding an agent, contracts, capability limits, privilege separation,
already exists, but it is built for organisations with security staff
and infrastructure behind it, not for a single person running an
assistant on their own machine. \artemis{}'s workspace model is best
read as a plain-language version of the same idea: a folder a user
selects rather than a sandbox a system administrator defines, and a
readable preview a person reads before an action runs rather than a
machine-checked contract. Where the trace-based framework has no way
to undo an action, \artemis{} adds exactly that as a deliberate design
choice, which is one of the more direct matches in this review. But
this needs to be stated plainly rather than smoothed over: \artemis{}
does not solve approval fatigue. It relies on the same human-in-the-loop
approval mechanism the security survey singles out as unresolved, and
nothing in the current design distinguishes a routine approval from a
risky one or reduces prompts as trust builds over repeated use. This
is a limitation the project inherits by choosing this approach at
all, and it is something the evaluation should test directly rather
than something the bounded workspace can be assumed to fix.

%% ------------------------------------------------------------------
\section{Agentic Reasoning, Planning, and Undo}
%% ------------------------------------------------------------------

A substantial body of work has moved language models beyond producing
text and toward reasoning about a goal, planning steps, and acting on
them. Reflexion lets an agent evaluate its own past attempts, write
verbal reflections, and store them as episodic memory to improve later
attempts, without retraining the model, and this improved performance
on decision-making, reasoning, and coding tasks
\autocite{shinn2023reflexion}. Toolformer takes a related but narrower
approach, teaching a model in a largely self-supervised way when and
how to call an external tool such as a calculator or a search engine
\autocite{schick2023toolformer}. Generative Agents extends the
memory-and-reflection idea further, showing that a population of
agents can behave coherently over simulated days by storing,
retrieving, and reflecting on experience, though the same paper also
reports failures such as incorrect retrieval and fabricated memory
details \autocite{park2023generative}. ReAct interleaves reasoning,
acting, and observing in a single loop and shows this combination
improves both performance and interpretability, while also being
prone to repetitive action loops and poor recovery when an action
returns an unhelpful result \autocite{yao2023react}. Tree of Thoughts
goes further still, letting a model search across multiple possible
reasoning paths and backtrack from unpromising ones, which produced
large gains on puzzle-style tasks but at a real computational cost
\autocite{yao2023tot}.

What unites this line of work is also its shared limitation. Each
system is evaluated on an individual task in a specific or simulated
environment: Reflexion on its own past attempts, Toolformer on
isolated tool calls it cannot easily chain together, Generative Agents
in a sandbox world, ReAct and Tree of Thoughts on benchmark tasks.
None of them is built to maintain a bounded workspace context for a
real user across ordinary, interrupted work, and none require the
kind of user approval that would make consequential actions
reversible. \artemis{} is positioned directly in that gap. It borrows
the underlying idea that reasoning, tool use, and reflective memory
can be combined into working agent behaviour, but it points that
behaviour at a user's actual files, notes, and stated intentions
instead of at the agent's own history or a simulated environment, and
it wraps the resulting multi-step plans in a preview-and-approval step
with an undo option, which none of these systems provide. The failure
modes reported across this literature, particularly fabricated memory
and repetitive reasoning loops, are also a direct warning for the
team: they describe exactly the kind of error the resumption and
synthesis features need to guard against, and they argue for keeping
the number of unsupervised steps between user checkpoints small rather
than assuming a longer autonomous chain is safe.

%% ------------------------------------------------------------------
\section{Theme Four: Voice Interaction Parity}
%% ------------------------------------------------------------------

\subsection{Overview of the Theme}

This theme concerns whether people can carry out practical computer work
by voice as well as they can by typing, and how voice interaction
should be evaluated. The reading covers two systematic reviews of voice
assistant usability, a study of speech as a real-time signal of user
experience, a study of speech-based waiting indicators, and a study of
conversational fillers for masking response delay.

\subsection{Deshmukh and Chalmeta, User Experience and Usability of
Voice User Interfaces (Information, 2024)}

\noindent Reference: \autocite{deshmukh2024vui}.

\subsubsection{What the Paper Found}

This systematic literature review conducts a bibliometric analysis of
user experience with voice user interfaces and proposes a
classification framework of six research categories. Existing research
mainly focuses on how users interact with voice assistants and what
factors affect their experience. The review identifies voice-related
factors such as voice gender, accent, personality, and style;
user-related factors such as gender, personality, experience, and
language background; task-related factors; conversational style; and
human-likeness or anthropomorphism. It also notes that voice-assistant
research measures trust, likeability, acceptance, perceived
intelligence, perceived humanness, social presence, and mental workload
in addition to the traditional effectiveness, efficiency, and
satisfaction of ISO 9241-11.

\subsubsection{Research Gap Identified or Left Open}

The review finds that existing research studies factors separately and
that the relationships between many variables and usability measures
are still incomplete. Some relationships are well studied while others
have very few studies. Traditional usability frameworks do not fully
capture voice-specific concerns such as trust, attitude, machine voice
and human-likeness, and cognitive load. Many studies use controlled
environments, and the review calls for more research in natural
settings. It also notes that voice as a complete interaction method for
multi-step desktop work is under-researched compared with single-task
voice use.

\subsubsection{What the Authors Were Trying to Do}

The authors were trying to organise a fragmented field, map its main
research streams, and identify future research directions within each
category.

\subsection{Dutsinma et al., A Systematic Review of Voice Assistant
Usability, ISO 9241-11 Approach (SN Computer Science, 2022)}

\noindent Reference: \autocite{dutsinma2022vareview}.

\subsubsection{What the Paper Found}

This review reads existing studies through the ISO 9241-11 lens of
effectiveness, efficiency, and satisfaction. It finds that voice
assistant research adds measures beyond that framework, including
trust, attitude, machine voice quality, human-likeness, and cognitive
load. It confirms that the same voice assistant may not give every user
the same experience, and that task type and context affect usability.

\subsubsection{Research Gap Identified or Left Open}

ISO 9241-11 alone does not explain all aspects of modern voice
assistant usability. Many aspects and relationships remain
insufficiently researched. Current usability measures do not completely
capture the needs of modern voice assistants, and there is a need for
broader evaluation of voice interaction in realistic, multi-step tasks,
especially when comparing voice with another interaction method such as
typing.

\subsubsection{What the Authors Were Trying to Do}

The authors were trying to apply a recognised usability standard
systematically to the voice-assistant literature and to show where that
standard is and is not sufficient.

\subsection{Ma et al., Beyond Words (CHI 2026)}

\noindent Reference: \autocite{ma2026beyondwords}.

\subsubsection{What the Paper Found}

This paper studies speech as a source of information about the user's
state during an interaction, not only after it. Human speech carries
vocal information such as pitch, speech rate, pauses, intensity,
jitter, shimmer, voice quality, filled pauses, repetitions, and
self-repairs. Prior work has connected these characteristics with
mental effort, frustration, engagement, uncertainty, and vocal tension,
and with user-experience concepts such as satisfaction, trust, and
attractiveness. The paper investigates whether speech features can be
used with machine-learning methods to distinguish positive, neutral,
and negative user experience automatically, in a within-subjects study
with participants.

\subsubsection{Research Gap Identified or Left Open}

Voice-interface evaluation still mainly depends on post-hoc
questionnaires and general performance measures, which ask users to
recall their experience after the interaction has finished and are
subject to recency bias, politeness, and limited self-awareness. There
is little evaluation that captures in-situ, turn-by-turn speech during
actual tasks. It is not clear which speech features reliably indicate
which user-experience dimensions, and whether system behaviour and
natural speech are sufficiently connected during real-time interaction.

\subsubsection{What the Authors Were Trying to Do}

The authors were trying to measure user experience from the user's
voice while the interaction is happening, by linking specific speech
features to user-experience levels through machine learning.

\subsection{Zhang et al., Human Speech Indicators for Waiting Experience
(CHI 2025)}

\noindent Reference: \autocite{speech2025waiting}.

\subsubsection{What the Paper Found}

The paper studies how human speech can be designed as a
waiting or loading indicator in voice user interfaces, where visual
indicators such as progress bars do not apply. Focus groups with
experienced voice-assistant users identified explanation and humour as
the two most important design requirements. The authors then tested
different speech indicators during short and long loading times. The
appropriate design depends on the length of the wait. For long waits, a
speech indicator with both high explanation and high humour performed
best overall. For short waits, either explanation or humour alone could
be effective, because too much information during a short wait can add
cognitive load.

\subsubsection{Research Gap Identified or Left Open}

Previous waiting-experience research focused on graphical interfaces
rather than voice interfaces, where delays are relatively invisible to
the user. There was not enough research on how different types of human
speech could be designed for the waiting period, and in particular the
combination of explanation and humour in speech indicators for voice
interfaces had not been properly investigated.

\subsubsection{What the Authors Were Trying to Do}

The authors were trying to find out how to fill the silent waiting
period in a voice interface with speech that improves the experience
rather than adding to the user's load.

\subsection{Wang et al., Mitigating Response Delays in Free-Form
Conversations (CUI 2025)}

\noindent Reference: \autocite{fillers2025mitigating}.

\subsubsection{What the Paper Found}

The authors studied response latency in free-form conversations with
language-model-powered virtual agents. They tested three response
delays of 1.5, 4, and 6.5 seconds and three conditions: no filler, an
artificial wait indicator, and a natural conversational filler. Longer
delays hurt users' perception of the agents. Natural conversational
fillers improved perceived response time, particularly at 4 and 6.5
seconds. Artificial wait indicators, consisting of a loading icon and
sound, did not significantly improve perceived response time or broader
user-experience measures compared with having no filler.

\subsubsection{Research Gap Identified or Left Open}

Earlier studies of response delay commonly used scripted questions and
answers, Wizard-of-Oz setups, or pre-recorded conversations rather than
fully interactive, free-form conversations with language-model agents,
which makes their findings hard to generalise. Simply showing the user
that the system is processing does not necessarily solve the
user-experience problem. Even natural fillers did not significantly
improve all broader aspects of user experience such as engagement, good
impression, competence, discomfort, or willingness to interact again.
Response delays above 4 seconds significantly degraded the experience,
and technical improvements cannot fully remove the problem because
network instability, hardware limits, and increasingly complex
reasoning continue to introduce latency.

\subsubsection{What the Authors Were Trying to Do}

The authors were trying to test whether conversational fillers reduce
the negative effect of response delay in realistic, free-form
interaction with a language-model agent, and to find the delay
threshold beyond which the experience degrades.

%% ------------------------------------------------------------------
\section{Theme Five: Code Assistance with Human Review}
%% ------------------------------------------------------------------

\subsection{Overview of the Theme}

This theme concerns artificial intelligence support for programming
that keeps a person in control of the decisions. The reading covers a
human-in-the-loop software development agent deployed in industry, a
token-level human-in-the-loop decoding interface, a study of when to
show a code suggestion, a study of expert steering strategies, and a
study of how reliably language models judge whether code meets a
requirement.

\subsection{Takerngsaksiri et al., Human-In-the-Loop Software
Development Agents (ICSE-SEIP 2025)}

\noindent Reference: \autocite{takerngsaksiri2025hula}.

\subsubsection{What the Paper Found}

The HULA framework has three agents: an artificial intelligence planner
agent, an artificial intelligence coding agent, and a human agent. The
planner identifies relevant files and creates a coding plan, the coding
agent generates code changes, and the human software engineer reviews,
modifies, and approves the results. HULA was integrated into Atlassian
JIRA and evaluated on both benchmark data and real issues. Practitioners
approved 433 of 527 generated plans, an 82 percent plan approval rate.
Of 376 issues for which code was generated, 95 resulted in raised pull
requests and 56 of those were merged. Practitioners reported that HULA
could reduce development time and effort, particularly for
straightforward tasks and for starting a coding plan.

\subsubsection{Research Gap Identified or Left Open}

Earlier systems were mostly evaluated on historical benchmark datasets
and open-source projects, which makes it hard to know how well they
work in real enterprise environments; previous work rarely included
human feedback at every stage; and many systems had not been deployed
in practice. HULA addresses these within software development but still
finds that agent performance depends heavily on the amount and quality
of input context, that generating consistently correct code without
human involvement remains a problem, and that passing unit tests does
not prove code is correct. HULA's workflow is specific to a JIRA issue,
a source-code repository, and pull requests, so it is an example of an
agent helping a human within one professional workflow rather than a
general desktop assistant.

\subsubsection{What the Authors Were Trying to Do}

The authors were trying to build and deploy a software development agent
that keeps a human in the loop at planning and coding stages, and to
report how real practitioners perceive its benefits and problems.

\subsection{Anaya Gonzalez et al., HiLDe (VL/HCC 2025)}

\noindent Reference: \autocite{gonzalez2025hilde}.

\subsubsection{What the Paper Found}

HiLDe introduces human-in-the-loop decoding. It highlights important
decision points in generated code, shows local alternatives that the
model considered, and explains how those alternatives differ. The user
selects a preferred alternative at the token level, and HiLDe
regenerates the following code according to that choice. In a study
with 18 participants, users generated 31 percent fewer vulnerabilities
on average with HiLDe than with a baseline assistant, and intentionally
repaired more vulnerabilities.

\subsubsection{Research Gap Identified or Left Open}

Existing assistants concentrate on code generation and provide global
alternatives that may be similar to one another, so users often have to
reverse-engineer the model's choices. Simply showing uncertainty is not
enough. HiLDe itself is still focused on code generation within a
programming environment, implemented as a Visual Studio Code extension
oriented to security-related tasks. The authors state that adapting it
to other priorities such as efficiency, maintainability, and personal
style is future work. HiLDe can also make unwanted changes to
downstream code when a user selects a local alternative, its
regeneration can overwrite previous user edits, and some participants
found the number of alternatives and the token highlighting
overwhelming.

\subsubsection{What the Authors Were Trying to Do}

The authors were trying to make interaction with an artificial
intelligence programming assistant more intentional, by surfacing
meaningful decisions and letting the user steer them at fine grain
rather than accepting the first suggestion.

\subsection{Mozannar et al., When to Show a Suggestion? (AAAI 2024)}

\noindent Reference: \autocite{mozannar2024whento}.

\subsubsection{What the Paper Found}

The paper presents an artificial intelligence-assisted programming
system built on GitHub Copilot that decides when to show or hide a code
suggestion. It uses programmer accept and reject behaviour as feedback,
learned from telemetry of real programming sessions covering 535
programmers and more than 168,000 shown suggestions. The main
contribution is conditional suggestion display from human feedback. A
machine-learning model predicts the probability that a programmer will
accept a suggestion; if acceptance is unlikely, the suggestion is
hidden. A two-stage design can decide using only code and prior
interaction information without generating the suggestion, which can
reduce unnecessary generation and latency. Including a model of the
programmer's latent state raised acceptance prediction accuracy from
61.9 percent to 83.6 percent.

\subsubsection{Research Gap Identified or Left Open}

The paper focuses on one specific interaction, deciding when to show or
hide a suggestion, and states that it does not tackle which suggestion
to display among a candidate set. The system relies on telemetry, which
cannot capture what the programmer is doing or thinking between two
recorded events, called the programmer's latent state. The evaluation
is mainly retrospective, using previously collected interaction data to
estimate what would have happened; the authors say a user study is
needed to verify whether the method makes programmers more productive,
and warn that hiding useful suggestions or relying too heavily on
acceptance rate can make the experience worse. Optimising only for
acceptance can favour very short or single-token suggestions rather
than complete useful code.

\subsubsection{What the Authors Were Trying to Do}

The authors were trying to reduce unnecessary suggestions, verification
effort, and latency in an artificial intelligence coding assistant, by
learning from accept and reject behaviour when a suggestion should be
displayed at all.

\subsection{Chidambaram et al., Socratic Human Feedback (Findings of
EMNLP 2024)}

\noindent Reference: \autocite{chidambaram2024sohf}.

\subsubsection{What the Paper Found}

The paper studies human steering of code-generating language models.
Eight expert programmers interacted with GPT-4, Gemini Ultra, and
Claude 3.5 Sonnet on difficult competition-level problems, using
multiple conversational turns to guide the models when their initial
solutions failed. The study identified steering strategies including
test reiteration, new test definition, revising unit tests, pointing
out a specific programming error, addressing code inefficiency,
requirement reiteration, requirement clarification, approach
re-orientation, and specific approach instruction. Humans could
successfully guide a model through several rounds of feedback rather
than expecting a correct solution in one attempt.

\subsubsection{Research Gap Identified or Left Open}

Previous work demonstrated human-model collaboration but there was
limited understanding of what specific feedback strategies expert users
actually use to steer a model toward a solution. This paper studies and
categorises those strategies rather than building them into a system.
It focuses on competition-level problems, and it is not yet clear
whether novice programmers would achieve the same success as experts,
since strategies such as pointing out a specific error and re-orienting
the approach may require expert knowledge. The dataset is relatively
small. Future work could use the identified strategies to design
interfaces that recommend follow-up prompts, ask clarifying questions,
or provide prompt templates, and could adapt steering strategies
dynamically to the conversation.

\subsubsection{What the Authors Were Trying to Do}

The authors were trying to observe real conversations between expert
programmers and code-generating models and to build a taxonomy of the
steering strategies those experts use.

\subsection{Jin and Chen, Are LLMs Reliable Code Reviewers? Systematic
Overcorrection in Requirement Conformance Judgement (Automated Software
Engineering, 2026)}

\noindent Reference: \autocite{jin2026reliable}.

\subsubsection{What the Paper Found}

The paper investigates whether language models can reliably verify that
code satisfies a natural-language requirement, especially when no test
cases or formal specifications are available. It finds that models make
two kinds of mistake: false rejection, where correct code is judged
incorrect, and false acceptance, where buggy code is judged correct. A
particularly important finding is an over-correction bias: instead of
checking whether the code meets the requirement, models invent
requirements, worry about unspecified edge cases, assume different
input types, imagine runtime problems, or claim the algorithm is wrong
when the implementation is actually correct. The largest category of
false-rejection explanations was logic error, followed by added
requirement, boundary error, and misread specification; together these
four categories represented 87.2 percent of false-negative
explanations. Increasing prompt complexity, for example asking the
model to judge, explain, and propose a fix, reduced false acceptance in
some situations but greatly increased false rejection. Explanations
could also contradict the model's own final judgement. The models were
better at identifying the observable symptom of a bug than the
underlying bug type.

\subsubsection{Research Gap Identified or Left Open}

Existing research shows that language models have strong
code-generation and code-understanding abilities, but much less is
known about whether their judgements themselves are reliable. Asking a
model for more reasoning does not necessarily make its reasoning more
trustworthy, and a longer explanation should not be treated as evidence
that a decision is correct.

\subsubsection{What the Authors Were Trying to Do}

The authors were trying to measure the reliability of language-model
judgements about requirement conformance, to build a taxonomy of
false-rejection reasons, and to test how prompt design affects
misjudgement rates.

%% ------------------------------------------------------------------
\section{Hallucination in Artificial Intelligence Summarisation}
%% ------------------------------------------------------------------

A separate concern applies to \artemis{}'s note and citation synthesis
feature: whether a generated summary stays faithful to the material it
is drawn from. The foundational study in this area ran a large human
evaluation across several summarisation systems and found substantial
hallucinated content in every one of them, splitting errors into
intrinsic cases, where the model distorts something from the source,
and extrinsic cases, where it adds something not present at all; it
also found that the standard ROUGE metric mostly failed to catch these
errors, while textual entailment between source and summary
correlated much better with human judgement
\autocite{maynez2020faithfulness}. This paper only measured the
problem; it did not build a detector.

Later work tried to close that gap directly. FactCC is a trained
classifier that predicts whether a summary sentence is consistent with
its source and can point to the specific span responsible for an
inconsistency, though it was trained on synthetic, rule-based errors,
which leaves some doubt about whether it generalises to the subtler
mistakes real generation systems make \autocite{kryscinski2020evaluating}.
A follow-up study argued that treating faithfulness as a single
yes-or-no judgement is too coarse, and built a finer typology of error
types together with a benchmark showing that no single automatic
metric, including FactCC, catches every kind of mistake reliably
\autocite{pagnoni2021understanding}. The most recent and arguably most
relevant paper in this theme tested actual large language models,
rather than older specialised classifiers, on the same underlying task
and found that most performed close to chance, with even GPT-4
landing several points below estimated human accuracy
\autocite{laban2023summedits}.

Read together, these findings say that hallucination in generated
summaries is common, that popular automatic metrics do not reliably
catch it, that no single detector covers every error type, and that
even strong general-purpose models are still poor judges of whether
their own output is faithful to a source. \artemis{}'s synthesis
feature narrows its exposure somewhat by working from a small,
explicitly user-selected set of documents rather than open-domain
summarisation, but a narrower scope is not the same as a solved
problem, and the current design has no faithfulness check of any
kind. The literature points toward a concrete, low-effort mitigation,
an entailment-style or FactCC-style consistency check applied after
generation, while also warning that such a check would not be a
complete fix on its own, and that asking the model to judge its own
citations is not a reliable substitute given how poorly even frontier
models perform at that task.

%% ------------------------------------------------------------------
\section{Retrieval-Augmented Generation Citation Accuracy}
%% ------------------------------------------------------------------

A closely related question is whether a citation attached to a
generated statement means anything at all, which matters directly for
the correctness of \artemis{}'s synthesised notes. One line of work
supplies a formal answer to what a good citation should look like,
proposing that a statement counts as properly attributed to a source
only if a person would agree that ``according to this source, this
statement is true,'' and building a human annotation procedure to
apply that test consistently across several generation tasks
\autocite{rashkin2023measuring}. This gives the field a definition,
but it depends on manual judgement and predates the current generation
of systems that attach citations automatically.

Later work measured how far short of that standard real systems fall.
A benchmark built specifically to test end-to-end citation generation
found that even the strongest models failed to fully support their
claims with citations roughly half the time on one of its test sets
\autocite{gao2023enabling}. An audit of four real, deployed generative
search products found much the same pattern under live conditions:
only about half of generated sentences were fully backed by their
citations, and roughly a quarter of the citations attached to a claim
did not actually support it \autocite{liu2023evaluating}.

The pattern across this theme is consistent: a clear definition of
what a trustworthy citation requires, evidence from a controlled
benchmark that current systems fall well short of it, and confirmation
from an audit of commercial products that the same gap shows up in
practice, including in systems with dedicated engineering behind their
citation pipelines. \artemis{} sidesteps part of this problem by
having the user pre-select the documents it draws from, which removes
the retrieval side of the difficulty these papers study. It does
nothing, however, about the generation side: nothing in the current
design checks that every claim in a synthesised note actually carries
a citation, or that each citation genuinely supports the statement it
is attached to. Given how often this fails even in polished, widely
used products, this should be treated as an open limitation of the
current design rather than something the narrower, user-selected
document set already resolves.

%% ==================================================================
\chapter{Critical Analysis and Takeaways for \artemis{}}
\label{ch:analysis}
%% ==================================================================

This chapter revisits every paper from Chapter~\ref{ch:raw} and adds the
element that the raw review deliberately left out: what the team should
take from it. For each paper the takeaway states what \artemis{} should
learn, whether \artemis{} addresses the gap the paper raises, and
whether that treatment is full or partial. Each theme closes with a
short synthesis.

%% ------------------------------------------------------------------
\section{Theme One: Task Interruption and Resumption}
%% ------------------------------------------------------------------

\subsection{Iqbal and Horvitz (CHI 2007)}

\subsubsection{Takeaway for \artemis{}}

\artemis{} turns this descriptive finding into a durable feature. The
\enquote{pick up where you left off} feature is a stable substitute for
the physical cue that Iqbal and Horvitz studied, because it surfaces the
equivalent information, what was recently active, what was saved, and
what the user said was next, as a persistent list rather than as a
window that happens to be open on a screen. This is a partial fill of
the gap. \artemis{} operationalises the finding but has not run the kind
of field validation the original study did, which is what the team's
field-evaluation objective is set up to do.

\subsection{Hu and Lee, ScreenTrack (CHI 2020)}

\subsubsection{Takeaway for \artemis{}}

\artemis{} avoids ScreenTrack's main weakness structurally rather than
trying to fix it. The resumption feature does not use screen recording
or visual scrubbing at all. It works from the file system and the
user's own saved notes, so it does not depend on whether two screens
look alike, and it matches on file name, recency, and explicit text
notes. It also sidesteps the \enquote{search a long video} problem,
because a workspace has a bounded scope by design, so there is a short
list to show rather than an open-ended timeline. The team should keep
the workspace boundary tight for exactly this reason.

\subsection{Synthesis for Theme One}

The two papers describe a progression: a field diagnosis with no tool,
followed by a tool that works only for recent, visually distinct
activity. \artemis{} sits deliberately to one side of this progression.
It does not attempt a general memory system. It substitutes a bounded
workspace, a short recency-and-notes list, and a visible editable
knowledge panel. The open work for the team is evaluation in a
realistic setting, since none of the resumption benefits have been
tested outside the design itself.

%% ------------------------------------------------------------------
\section{Theme Two: Bounded Artificial Intelligence Agents and Trust}
%% ------------------------------------------------------------------

\subsection{Paduraru et al., Trace-Based Assurance (2026)}

\subsubsection{Takeaway for \artemis{}}

This is one of the clearest direct matches in the whole review. The
framework's biggest blind spot, no way to reverse an action once it
executes, is exactly what \artemis{}'s undo step is built to cover.
Where this framework mediates actions using machine-checkable contracts
before they run, \artemis{} mediates them with a plain-language preview
that a person reads before they run, and then adds the reversal option
this framework never has. It is a partial fill in the other direction:
\artemis{} has no deterministic replay and no automatic contract
verification, so it depends entirely on a person looking at a preview
and deciding, which is a weaker guarantee than a machine-checked
contract but a more usable one for someone without an engineering
background.

\subsection{Kim et al., Attack and Defense Landscape (2026)}

\subsubsection{Takeaway for \artemis{}}

This one needs to be stated plainly rather than smoothed over.
\artemis{} does not solve approval fatigue. This paper identifies a
real, documented weakness in the exact approach \artemis{} depends on,
stopping for approval before every consequential action. The design
currently has no mechanism to prevent that fatigue, nothing that
distinguishes a routine low-stakes approval from a genuinely risky one,
and nothing that reduces prompts as trust builds. This is a limitation
\artemis{} inherits from relying on human-in-the-loop approval at all,
not something its bounded workspace design removes. The team's
evaluation should test directly whether approval fatigue shows up over
repeated sessions, rather than assuming the design avoids it.

\subsection{Synthesis for Theme Two}

Across the two papers, the safe-agent field has strong ideas, contract
mediation, deterministic replay, capability limits, and privilege
separation, but they are built for developers and organisations with
security staff. \artemis{} takes the principles and makes them operable
by one person: a folder instead of a sandbox definition, a readable
panel instead of a provenance graph, a preview and undo instead of a
contract engine. The honest gap is approval fatigue, which \artemis{}
inherits and must measure rather than claim to have designed away.

%% ------------------------------------------------------------------
\section{Theme Three: Agentic Reasoning, Planning, and Undo}
%% ------------------------------------------------------------------

\subsection{Shinn et al., Reflexion (2023)}

\subsubsection{Takeaway for \artemis{}}

\artemis{} applies agentic reasoning in a bounded desktop workspace.
Instead of focusing only on learning from failed attempts, it uses
workspace context such as recently changed files, saved notes, and the
user's intended next action to help reconstruct interrupted work. It
combines multi-step planning with explicit workspace boundaries and
reversible actions. The lesson from Reflexion is that episodic,
written memory is a workable substrate for agent behaviour; \artemis{}
uses a similar idea but points it at human work context rather than at
the agent's own past errors.

\subsection{Schick et al., Toolformer (2023)}

\subsubsection{Takeaway for \artemis{}}

\artemis{} extends tool-using language models toward chained,
multi-step execution, which Toolformer does not fully support. A request
such as reading documents, preparing a summary, and emailing it is a
sequence where each step depends on the previous one. \artemis{} plans
that sequence and keeps it bounded and reviewable through the approval
mechanism. The takeaway is that tool use alone is not enough for a real
assistant; the planning layer on top is where \artemis{} adds value.

\subsection{Park et al., Generative Agents (2023)}

\subsubsection{Takeaway for \artemis{}}

\artemis{} applies the memory, reflection, and planning idea to a
practical desktop workspace instead of a simulated world. Given a
request, the agent determines what information is needed, retrieves
relevant workspace context, selects tools, and plans a sequence of
actions. Unlike a fully autonomous agent it requires user approval
before actions that modify files or leave the workspace, and it
provides preview, audit, and undo. The reported failures in this paper,
incorrect memory retrieval and fabricated memory details, are a direct
warning: they are the same failure modes the team must guard against in
the resumption and synthesis features.

\subsection{Yao et al., ReAct (2023)}

\subsubsection{Takeaway for \artemis{}}

\artemis{} adopts the interleaved reason, act, observe loop but places
it inside a bounded workspace with user approval at the points where
actions leave that workspace. The known weaknesses of ReAct, reasoning
errors, repetitive loops, and poor recovery from uninformative
observations, argue for keeping the action space small and the number
of autonomous steps limited before returning to the user. The team
should treat the approval checkpoints as loop breakers as well as trust
mechanisms.

\subsection{Yao et al., Tree of Thoughts (2023)}

\subsubsection{Takeaway for \artemis{}}

Tree of Thoughts establishes that search, evaluation, and backtracking
improve reasoning, but its evaluation is on puzzle-like tasks. \artemis{}
does not need full tree search; its tasks are more like short
dependency chains than large search spaces. The takeaway is that a
lightweight planning step, with the ability to revise a plan when a
step produces a poor result, is enough for the \artemis{} task profile,
and heavier search would add computation without a matching benefit.

\subsection{Synthesis for Theme Three}

The five papers establish that language models can reason, plan, use
tools, reflect, and search, and that reasoning alone is not enough for a
real assistant. The common gap is that this work is evaluated on
individual tasks in specific or simulated environments, not on
maintaining a bounded workspace context for a user across everyday
activity. \artemis{} is positioned exactly in that gap: it applies
multi-step agentic planning to a practical, bounded workspace, with
tool selection, context awareness, human approval, and reversible
actions. The failure modes reported across these papers, especially
fabricated memory and reasoning loops, define the risks the team's
evaluation must probe.

%% ------------------------------------------------------------------
\section{Theme Four: Voice Interaction Parity}
%% ------------------------------------------------------------------

\subsection{Deshmukh and Chalmeta (2024)}

\subsubsection{Takeaway for \artemis{}}

This review tells the team what has and has not been studied in voice
usability, and confirms that voice as a complete interaction method for
multi-step desktop work is under-researched. \artemis{} contributes here
by treating speech as an alternative interaction mode over the same
bounded feature set, and by evaluating it against typed interaction on
ease of use and task completion time, with cognitive workload and trust
in the broader evaluation. This is a partial fill: \artemis{} studies
voice in a real desktop productivity setting rather than as a
question-and-answer interface, but it does not resolve the fragmented
factor relationships the review describes.

\subsection{Dutsinma et al. (2022)}

\subsubsection{Takeaway for \artemis{}}

The review's central point, that ISO 9241-11's effectiveness,
efficiency, and satisfaction do not fully capture voice usability, tells
the team which extra measures to include. The \artemis{} voice
evaluation should therefore report not only task completion time and
success but also cognitive workload and trust, which this review
identifies as voice-specific concerns. The safer framing for the team's
own writing is that the reviewed literature does not sufficiently
address voice interaction parity across complete, multi-step desktop
workflows, rather than any claim that voice and typing have never been
compared.

\subsection{Ma et al., Beyond Words (2026)}

\subsubsection{Takeaway for \artemis{}}

This paper is relevant to how the team measures the voice mode, not to
what \artemis{} builds. It shows that speech carries a real-time signal
of user experience that post-hoc questionnaires miss. If the team can
record turn-by-turn audio during the voice-versus-typing study, speech
features such as pause duration and disfluency rate could supplement
questionnaire measures of cognitive load and frustration. \artemis{}
does not currently plan this, so it is an optional evaluation
enhancement rather than a design requirement.

\subsection{Zhang et al., Waiting Indicators (2025)}

\subsubsection{Takeaway for \artemis{}}

\artemis{}'s voice mode will have waiting periods while the agent plans
and executes. This paper gives concrete design guidance: for a long
wait, a spoken indicator that both explains what is happening and adds
light humour works best; for a short wait, explanation or humour alone
is enough, and too much information adds load. The team should design
the voice mode's waiting behaviour along these lines rather than using a
silent pause or a generic tone.

\subsection{Wang et al., Response Delays (2025)}

\subsubsection{Takeaway for \artemis{}}

This paper sets a latency budget for the voice mode. Response delays
above 4 seconds significantly degrade the experience, and natural
conversational fillers help at 4 and 6.5 seconds while artificial
loading sounds do not. The team should target sub-4-second responses
where possible, use natural spoken fillers rather than earcons when a
delay is unavoidable, and accept that filler cannot fully compensate for
a slow multi-step plan. This reinforces keeping the number of
autonomous steps between user touchpoints small.

\subsection{Synthesis for Theme Four}

The two reviews establish that voice usability is studied in fragments
and that standard usability metrics miss voice-specific concerns such as
trust and cognitive load. The three empirical papers give the team
usable numbers and design rules: a 4-second latency threshold, natural
fillers over artificial ones, explanation-plus-humour for long waits,
and speech as an optional real-time experience signal. \artemis{}'s
contribution is to evaluate voice as a full interaction mode over
practical multi-step tasks, comparing it against typing on ease of use,
task time, cognitive workload, and trust, which the literature agrees is
under-explored.

%% ------------------------------------------------------------------
\section{Theme Five: Code Assistance with Human Review}
%% ------------------------------------------------------------------

\subsection{Takerngsaksiri et al., HULA (2025)}

\subsubsection{Takeaway for \artemis{}}

HULA is evidence that a plan, generate, human-approve loop works in
practice, with an 82 percent plan approval rate in a real deployment.
\artemis{}'s code assistant uses the same shape but is narrower: it
reads workspace source, answers questions, explains errors, and
proposes a fix as a diff that the user approves, and it does not open
pull requests or act on a whole repository. The lesson is that the
human-approval loop is not a compromise; it is a validated interaction
model. HULA's finding that output quality depends heavily on input
context also applies: \artemis{} should surface the relevant files and
context clearly before proposing a change.

\subsection{Anaya Gonzalez et al., HiLDe (2025)}

\subsubsection{Takeaway for \artemis{}}

HiLDe operates at a finer grain than \artemis{} needs, steering
generation token by token. \artemis{} keeps the interaction at the level
of a whole proposed change shown as a diff, which is coarser but
simpler and matches the non-expert user. The relevant caution from
HiLDe is that regeneration can silently overwrite earlier user edits and
that too many alternatives overwhelm users; \artemis{} should show one
proposed change at a time and never modify code the user did not
approve.

\subsection{Mozannar et al., When to Show a Suggestion? (2024)}

\subsubsection{Takeaway for \artemis{}}

This paper addresses a problem \artemis{} does not have. \artemis{}'s
code assistant is invoked deliberately by the user rather than offering
continuous inline suggestions, so there is no when-to-show decision to
make and no risk of optimising for acceptance rate at the expense of
useful content. The takeaway is that \artemis{}'s on-demand model
sidesteps a whole class of problems that inline assistants have, and the
team can note this as a deliberate design advantage.

\subsection{Chidambaram et al., Socratic Human Feedback (2024)}

\subsubsection{Takeaway for \artemis{}}

This paper's taxonomy of steering strategies, such as pointing out a
specific error, clarifying a requirement, and re-orienting the approach,
is useful input for \artemis{}'s code-assistant conversation design. The
team could use it to structure follow-up prompts or clarifying
questions when the user disagrees with a proposed fix. The paper's
caution that some strategies may need expert knowledge is relevant,
since \artemis{}'s target user may be a non-expert, so the assistant
should offer these steering moves rather than assume the user will
supply them.

\subsection{Jin and Chen, Are LLMs Reliable Code Reviewers? (2026)}

\subsubsection{Takeaway for \artemis{}}

This paper is the direct justification for \artemis{} never applying a
code change on its own. It adds a specific warning: models over-correct,
inventing requirements and flagging correct code as wrong, and asking
for a longer explanation makes this worse, not better. Increasing
prompt complexity reduced false acceptance in some cases but greatly
increased false rejection, and explanations could contradict the
model's own final judgement. \artemis{} should therefore present a
proposed fix as a suggestion to review, not as a defect report to
trust, and should be cautious about prompting its own model for
elaborate justifications. The paper's point that the models were better
at spotting a bug's symptom than its underlying cause also tells the
team to keep the user, not the model, as the final judge. This is a
limitation \artemis{} inherits from using a language model at all, and
the diff-plus-approval design is the mitigation.

\subsection{Synthesis for Theme Five}

The five papers converge on one point: language-model code judgement is
useful but not trustworthy on its own, with a documented over-correction
bias and judgements that a longer explanation does not make more
reliable. \artemis{}'s design answers this directly through an
on-demand, diff-based, user-approved code assistant that never writes or
runs code by itself. HULA and HiLDe show the human-approval loop is a
validated model, not a fallback. The steering-strategy taxonomy from
SoHF and the acceptance-versus-usefulness caution from the
when-to-show-a-suggestion work are concrete inputs for the assistant's
conversation design. The residual limitation, that the underlying model
can still be wrong or over-cautious, is inherited and is mitigated, not
removed, by keeping the user in control.

%% ------------------------------------------------------------------
\section{Theme Six: Hallucination in Artificial Intelligence
Summarisation}
%% ------------------------------------------------------------------

\subsection{Maynez et al. (2020)}

\subsubsection{Takeaway for \artemis{}}

The synthesis feature pulls together notes and citations from a small,
explicitly selected set of documents, which is a narrower and more
controlled task than open-domain abstractive summarisation. A narrower
scope reduces the surface area for hallucination but does not remove it.
This paper's finding, that hallucination appeared in every system
tested, is a real reason to expect the synthesis feature needs some kind
of faithfulness check, and the team should not treat the narrower scope
as making the problem solved. Entailment between source and output is
the check this paper points to.

\subsection{Kryscinski et al., FactCC (2020)}

\subsubsection{Takeaway for \artemis{}}

FactCC is the closest thing in the literature to a plug-in tool the
synthesis feature could actually use: a lightweight, already-built
consistency checker rather than something built from scratch. Nothing
in the current \artemis{} design uses a detector of this kind, so this
is a genuine open gap rather than something \artemis{} already handles.
If the team wants to reduce hallucination risk in the synthesis feature,
adopting a FactCC-style check after generation is a concrete,
literature-backed way to do it.

\subsection{Pagnoni et al., FRANK (2021)}

\subsubsection{Takeaway for \artemis{}}

This paper is a warning against relying on a single automatic metric to
prove the synthesis feature is faithful. If the team adds a consistency
check, FRANK's finding means one metric will not catch every kind of
error a citation-synthesis feature can produce. It does not hand
\artemis{} a gap to fill so much as it tells the team to expect any
single fix to be incomplete and to combine checks or keep a human in
the loop.

\subsection{Laban et al., SummEdits (2023)}

\subsubsection{Takeaway for \artemis{}}

This is the strongest reason to treat the synthesis feature's evaluation
as genuinely necessary rather than a formality. It shows that even
frontier models cannot reliably tell a faithful summary from an
unfaithful one, so if \artemis{} ever tried to have the model check its
own citations by asking \enquote{is this supported by the source}, that
approach is not reliable even with the best models. This is a current
limitation \artemis{} inherits, and the team's plan to measure synthesis
accuracy where source metadata is available is the right response.

\subsection{Synthesis for Theme Six}

The four papers trace the problem from diagnosis to detection to
meta-evaluation to a frontier-model test, and the conclusion is
consistent: hallucination is common, ROUGE does not catch it, no single
detector is uniformly reliable, and even the best current models are
weak at judging faithfulness. \artemis{} reduces exposure by narrowing
the task to a small, user-selected document set, but it currently has no
faithfulness check, so this is an open gap in the design. The
recommended direction is an entailment-style or FactCC-style check after
generation, combined with keeping the user in the loop, plus the planned
accuracy evaluation where source metadata allows it.

%% ------------------------------------------------------------------
\section{Theme Seven: Retrieval-Augmented Generation Citation Accuracy}
%% ------------------------------------------------------------------

\subsection{Rashkin et al., Attribution (2023)}

\subsubsection{Takeaway for \artemis{}}

The \enquote{according to this source, this statement is true} test is a
useful standard to hold the synthesis feature's output to: every claim
it pulls into a new file should pass that test against the document it
came from. This paper gives \artemis{} a standard to aim for, not a tool
to enforce it automatically, and nothing in the current design applies
this standard even informally. Adding a check of this kind, even a
manual review step during evaluation, would be a concrete improvement.

\subsection{Gao et al., ALCE (2023)}

\subsubsection{Takeaway for \artemis{}}

Having the user pre-select the documents removes the retrieval half of
the problem ALCE studies, since \artemis{} never has to decide which
documents to retrieve. But ALCE's finding that even strong systems fail
to fully cite their claims about half the time is about the generation
half, and a smaller document set does not fix that. Nothing in the
current design checks that every claim in a synthesised file carries a
citation, so citation completeness is an open gap that the workspace
scope does not resolve on its own.

\subsection{Liu et al., Verifiability in Generative Search (2023)}

\subsubsection{Takeaway for \artemis{}}

This is the clearest real-world comparison point for the synthesis
feature: mature, widely used products with dedicated citation
engineering still fail these checks often, with about a quarter of
citations not holding up and about half of sentences not fully
supported. That is a reason to expect the synthesis feature needs
deliberate citation verification, not an assumption that citing from a
small known set is inherently safe. The proposal narrows the retrieval
problem by having the user pre-select documents, but it does not
currently include the recall or precision checks this paper used to
audit real products, so this is an unaddressed gap in the current
design, not a partially solved one.

\subsection{Synthesis for Theme Seven}

The three papers give a definition of attribution, a benchmark showing
incomplete citation is common, and an audit showing deployed products
get citation precision and recall wrong at meaningful rates. \artemis{}
removes the retrieval part of the problem by having the user choose the
documents, but citation correctness and completeness in the generated
output remain unaddressed in the current design. The team should either
add a citation-verification step, at least for the evaluation, or be
explicit in the project report that this is a known limitation rather
than something the bounded workspace solves.

%% ==================================================================
\chapter{Cross-Cutting Observations}
%% ==================================================================

\section{Where \artemis{} Genuinely Fills a Gap}

\begin{enumerate}
  \item Applying multi-step agentic planning to a bounded, user-chosen
        desktop workspace, rather than to a simulated world, a single
        task, or an open codebase. The reasoning-and-planning papers
        all leave this specific combination open.
  \item Making agent memory visible and user-editable through the
        \enquote{what \artemis{} knows} panel, rather than running it as
        an opaque background process. The Generative Agents work reports
        fabricated and mis-retrieved memory, and the agent-security
        survey notes that persistence is usually an automatic process
        the user cannot inspect or control.
  \item Providing a reversal step after a previewed action. The
        trace-based assurance framework has no reversal mechanism at
        all, and this is one of the clearest direct matches in the
        review.
  \item Evaluating voice as a complete interaction mode over practical
        multi-step tasks, compared against typing on ease of use, task
        time, cognitive workload, and trust. Both voice reviews agree
        this is under-explored.
\end{enumerate}

\section{Where \artemis{} Only Partially Fills a Gap}

\begin{enumerate}
  \item Bounded capability scoping. \artemis{}'s workspace is a
        plain-language version of the capability limits and privilege
        separation the assurance and security papers describe, and it
        trades rigour for usability.
  \item Task resumption. \artemis{} operationalises the Iqbal and
        Horvitz finding but has not run comparable field validation.
  \item Code assistance. The human-approval loop is validated by HULA
        and HiLDe, but the underlying model can still be wrong or
        over-cautious, and approval mitigates rather than removes this.
\end{enumerate}

\section{Where \artemis{} Does Not Fill the Gap and Inherits a
Limitation}

\begin{enumerate}
  \item Approval fatigue. The security survey documents that frequent
        approval prompts cause decision fatigue and that non-expert
        users cannot always judge what they approve. \artemis{} relies
        on approval prompts and has no mechanism to vary them by risk
        or to reduce them as trust builds. The evaluation must test for
        this rather than assume the design avoids it.
  \item Faithfulness of synthesised text. Every summarisation paper
        shows hallucination is common and hard to detect, and even
        frontier models judge faithfulness poorly. \artemis{} narrows
        the task but has no faithfulness check in the current design.
  \item Citation correctness and completeness. The attribution and
        citation papers show that citing is unreliable even in polished
        deployed products. \artemis{} removes the retrieval half by
        pre-selection but does not verify that citations support their
        claims or that every claim is cited.
\end{enumerate}

\section{Implications for the Evaluation Plan}

\begin{enumerate}
  \item The field-style evaluation should measure task resumption
        benefit directly, since no reviewed work has validated this
        outside a design description.
  \item The evaluation should include a specific test for approval
        fatigue over repeated sessions.
  \item The voice study should use a 4-second latency threshold as a
        design target, natural spoken fillers rather than tones, and
        should report cognitive workload and trust alongside task time
        and success.
  \item The synthesis feature should be evaluated for factual
        faithfulness and citation accuracy wherever source metadata
        allows it, and any shortfall should be reported as a known
        limitation.
  \item Code-assistant evaluation should confirm that no change is ever
        applied without explicit user approval and should check for
        silent overwrites of prior user edits.
\end{enumerate}

%% ==================================================================
\printbibliography[heading=bibintoc,title={References}]
%% ==================================================================

\end{document}
